\PassOptionsToPackage{unicode}{hyperref}
\PassOptionsToPackage{hyphens}{url}
\documentclass[11pt,a4paper]{article}

% ── Codificação e língua ──────────────────────────────────────────────────────
\usepackage[T1]{fontenc}
\usepackage[utf8]{inputenc}
\usepackage[brazil]{babel}

% ── Tipografia ────────────────────────────────────────────────────────────────
\usepackage{lmodern}
\usepackage{microtype}

% ── Cores ─────────────────────────────────────────────────────────────────────
\usepackage{xcolor}
\definecolor{javaOrange}{HTML}{E76F00}
\definecolor{javaDeep}{HTML}{BF5700}
\definecolor{javaBlue}{HTML}{1565C0}
\definecolor{javaTeal}{HTML}{00695C}
\definecolor{javaAmber}{HTML}{B45309}
\definecolor{javaInk}{HTML}{2B2140}
\definecolor{javaCodeBg}{HTML}{FFF8F0}
\definecolor{javaRule}{HTML}{FFD7B5}
\definecolor{javaShade}{HTML}{FFF3E8}

% ── Caixas e código ───────────────────────────────────────────────────────────
\usepackage{tcolorbox}
\tcbuselibrary{skins,breakable,listings}
\usepackage{listings}
\usepackage{fvextra}

\lstdefinestyle{javastyle}{
  language=Java,
  basicstyle=\ttfamily\small,
  keywordstyle=\color{javaBlue}\bfseries,
  stringstyle=\color{javaDeep},
  commentstyle=\color{javaTeal}\itshape,
  numberstyle=\tiny\color{gray},
  breaklines=true,
  breakatwhitespace=false,
  showstringspaces=false,
  tabsize=4,
}

\newtcblisting{javacode}{
  listing only,
  listing style=javastyle,
  breakable,
  enhanced,
  colback=javaCodeBg,
  colframe=javaDeep,
  leftrule=4pt,
  rightrule=0.4pt,
  toprule=0.4pt,
  bottomrule=0.4pt,
  arc=2pt,
  fontupper=\small,
  top=4pt, bottom=4pt, left=6pt, right=6pt,
}

\newtcbox{\inlinecode}{
  on line,
  colback=javaCodeBg,
  colframe=javaRule,
  boxrule=0.5pt,
  arc=2pt,
  fontupper=\ttfamily\small\color{javaDeep},
  boxsep=1pt, left=2pt, right=2pt, top=1pt, bottom=1pt,
}

% ── Layout ────────────────────────────────────────────────────────────────────
\usepackage[margin=2.4cm, top=2.8cm, bottom=2.8cm]{geometry}

% ── Cabeçalho / rodapé ────────────────────────────────────────────────────────
\usepackage{fancyhdr}
\pagestyle{fancy}
\fancyhf{}
\renewcommand{\headrulewidth}{0.4pt}
\renewcommand{\headrule}{\hbox to\headwidth{\color{javaRule}\leaders\hrule height \headrulewidth\hfill}}
\fancyhead[L]{\small\sffamily\color{javaDeep} Java Programming MOOC}
\fancyhead[R]{\small\sffamily\color{javaDeep} Lição 1}
\fancyfoot[C]{\small\sffamily\color{gray}\thepage}

% ── Seções ────────────────────────────────────────────────────────────────────
\usepackage{titlesec}
\titleformat{\section}
  {\sffamily\Large\bfseries\color{javaOrange}}
  {\thesection}{0.7em}{}
  [\vspace{2pt}{\color{javaRule}\titlerule[1.2pt]}]
\titleformat{\subsection}
  {\sffamily\large\bfseries\color{javaDeep}}
  {\thesubsection}{0.6em}{}
\titleformat{\subsubsection}
  {\sffamily\normalsize\bfseries\color{javaTeal}}
  {\thesubsubsection}{0.5em}{}

% ── Listas ────────────────────────────────────────────────────────────────────
\usepackage{enumitem}
\setlist[itemize]{itemsep=2pt, topsep=4pt, parsep=0pt}
\setlist[enumerate]{itemsep=2pt, topsep=4pt, parsep=0pt}

% ── Tabelas ───────────────────────────────────────────────────────────────────
\usepackage{booktabs}
\usepackage{array}
\usepackage{colortbl}
\usepackage{longtable}

% ── Hiperlinks ────────────────────────────────────────────────────────────────
\usepackage{hyperref}
\hypersetup{
  colorlinks=true,
  linkcolor=javaBlue,
  urlcolor=javaBlue,
  citecolor=javaBlue,
  pdftitle={Programação em Java — Lição 1: Primeiros Passos},
  pdfauthor={Java Programming MOOC · Universidade de Helsinque},
  pdflang={pt-BR},
}

% ── Caixa de destaque (dica/atenção) ─────────────────────────────────────────
\newtcolorbox{destaque}[1][Dica]{
  enhanced, breakable,
  colback=javaShade,
  colframe=javaOrange,
  leftrule=4pt, rightrule=0.4pt, toprule=0.4pt, bottomrule=0.4pt,
  arc=2pt,
  fonttitle=\sffamily\bfseries\color{javaOrange},
  title=#1,
  top=4pt, bottom=4pt,
}

% ── Título personalizado ──────────────────────────────────────────────────────
\makeatletter
\newcommand{\subtitle}[1]{\gdef\@subtitle{#1}}
\renewcommand{\maketitle}{%
  \begin{center}
    {\color{javaRule}\rule{\linewidth}{1pt}}\\[6pt]
    {\Huge\sffamily\bfseries\color{javaOrange} \@title}\\[6pt]
    \ifx\@subtitle\undefined\else
      {\Large\sffamily\color{javaDeep} \@subtitle}\\[4pt]
    \fi
    {\small\sffamily\color{gray} \@author}\\[2pt]
    {\small\sffamily\color{gray} \@date}\\[6pt]
    {\color{javaRule}\rule{\linewidth}{1pt}}
  \end{center}
  \vspace{1em}
}
\makeatother

% ─────────────────────────────────────────────────────────────────────────────
\title{Programação em Java}
\subtitle{Lição 1: Primeiros Passos}
\author{Java Programming MOOC · Universidade de Helsinque — tradução para o português}
\date{2026}
% ─────────────────────────────────────────────────────────────────────────────

\begin{document}
\maketitle
\tableofcontents
\vspace{1.5em}

% ─────────────────────────────────────────────────────────────────────────────
\section{Objetivos de aprendizagem}

Ao final desta lição, você será capaz de:

\begin{itemize}
  \item Entender o que é um programa de computador
  \item Escrever programas que exibem texto na tela
  \item Ler dados digitados pelo usuário
  \item Criar e usar variáveis dos tipos \inlinecode{String}, \inlinecode{int}, \inlinecode{double} e \inlinecode{boolean}
  \item Realizar cálculos aritméticos básicos
  \item Usar condicionais (\inlinecode{if}, \inlinecode{else if}, \inlinecode{else})
\end{itemize}

% ─────────────────────────────────────────────────────────────────────────────
\section{O que é programação?}

Programação é a arte de escrever instruções que um computador possa executar. Essas instruções são chamadas de \textbf{código}, e um conjunto organizado delas forma um \textbf{programa}.

Os computadores são máquinas que executam instruções com extrema rapidez, mas são completamente literais: fazem exatamente o que o programador manda, nem mais, nem menos. Por isso, escrever código envolve ser muito preciso.

Java é uma das linguagens de programação mais utilizadas no mundo. É uma linguagem \textbf{compilada}, \textbf{fortemente tipada} e \textbf{orientada a objetos}. Neste curso, aprenderemos Java do zero.

% ─────────────────────────────────────────────────────────────────────────────
\section{Estrutura de um programa Java}

Todo programa Java tem a mesma estrutura básica:

\begin{javacode}
public class NomeDoPrograma {
    public static void main(String[] args) {
        // o código do programa vai aqui
    }
}
\end{javacode}

\begin{itemize}
  \item \inlinecode{public class NomeDoPrograma} define a \textbf{classe} do programa. O nome da classe deve coincidir com o nome do arquivo (ex.: \inlinecode{NomeDoPrograma.java}).
  \item \inlinecode{public static void main(String[] args)} é o \textbf{método principal} — o ponto de partida da execução.
  \item As chaves \inlinecode{\{ \}} delimitam blocos de código.
  \item O \inlinecode{//} inicia um comentário, que é ignorado pelo compilador.
\end{itemize}

% ─────────────────────────────────────────────────────────────────────────────
\section{Exibindo texto na tela}

O comando mais básico de Java é o \inlinecode{System.out.println}, que imprime uma linha de texto:

\begin{javacode}
public class OlaMundo {
    public static void main(String[] args) {
        System.out.println("Olá, mundo!");
    }
}
\end{javacode}

\textbf{Saída:}
\begin{javacode}
Olá, mundo!
\end{javacode}

Você pode imprimir várias linhas repetindo o comando:

\begin{javacode}
System.out.println("Primeira linha");
System.out.println("Segunda linha");
System.out.println("Terceira linha");
\end{javacode}

Para imprimir sem pular a linha, use \inlinecode{System.out.print}:

\begin{javacode}
System.out.print("Olá, ");
System.out.print("mundo!");
\end{javacode}

\begin{destaque}[Dica]
Na maioria das IDEs, você pode digitar \inlinecode{sout} e pressionar Tab para escrever \inlinecode{System.out.println()} automaticamente.
\end{destaque}

% ─────────────────────────────────────────────────────────────────────────────
\section{Lendo a entrada do usuário}

Para ler dados digitados pelo usuário, usamos a classe \inlinecode{Scanner}:

\begin{javacode}
import java.util.Scanner;

public class LeitorDeNome {
    public static void main(String[] args) {
        Scanner scanner = new Scanner(System.in);

        System.out.print("Qual é o seu nome? ");
        String nome = scanner.nextLine();

        System.out.println("Olá, " + nome + "!");
    }
}
\end{javacode}

O \inlinecode{import java.util.Scanner} torna a classe \inlinecode{Scanner} disponível. O \inlinecode{scanner.nextLine()} lê uma linha inteira de texto digitada pelo usuário.

% ─────────────────────────────────────────────────────────────────────────────
\section{Variáveis}

Uma \textbf{variável} é um espaço nomeado na memória onde guardamos um valor. Em Java, toda variável tem um \textbf{tipo} definido na criação.

\subsection{Tipos básicos}

\begin{center}
\begin{tabular}{>{\ttfamily}l l >{\ttfamily}l}
\toprule
\normalfont\textbf{Tipo} & \textbf{Descrição} & \normalfont\textbf{Exemplo} \\
\midrule
String  & Texto               & "Olá" \\
int     & Número inteiro      & 42    \\
double  & Número decimal      & 3.14  \\
boolean & Verdadeiro ou falso & true, false \\
\bottomrule
\end{tabular}
\end{center}

\subsection{Declarando variáveis}

\begin{javacode}
String nome = "Ana";
int idade = 25;
double altura = 1.68;
boolean estudante = true;

System.out.println(nome);      // Ana
System.out.println(idade);     // 25
System.out.println(altura);    // 1.68
System.out.println(estudante); // true
\end{javacode}

\subsection{Boas práticas de nomeação}

\begin{itemize}
  \item Use \textbf{camelCase}: primeira palavra em minúsculo, demais começam em maiúsculo.
  \item Escolha nomes significativos.
\end{itemize}

\begin{javacode}
// Ruim
double a = 3.14, b = 22.0, c = a * b * b;

// Bom
double pi = 3.14, raio = 22.0;
double areaDoCirculo = pi * raio * raio;
\end{javacode}

\subsection{Lendo valores de outros tipos}

\begin{javacode}
Scanner scanner = new Scanner(System.in);

System.out.print("Digite um número inteiro: ");
int numero = Integer.valueOf(scanner.nextLine());

System.out.print("Digite um número decimal: ");
double decimal = Double.valueOf(scanner.nextLine());
\end{javacode}

% ─────────────────────────────────────────────────────────────────────────────
\section{Calculando com números}

\begin{center}
\begin{tabular}{c l l l}
\toprule
\textbf{Operador} & \textbf{Operação} & \textbf{Exemplo} & \textbf{Resultado} \\
\midrule
\texttt{+} & Adição         & \texttt{3 + 4}  & \texttt{7}           \\
\texttt{-} & Subtração      & \texttt{10 - 3} & \texttt{7}           \\
\texttt{*} & Multiplicação  & \texttt{5 * 6}  & \texttt{30}          \\
\texttt{/} & Divisão        & \texttt{10 / 3} & \texttt{3} (inteira!) \\
\texttt{\%} & Resto         & \texttt{10 \% 3} & \texttt{1}          \\
\bottomrule
\end{tabular}
\end{center}

\subsection{Divisão inteira vs.\ divisão decimal}

Em Java, dividir dois \inlinecode{int} sempre dá resultado inteiro:

\begin{javacode}
System.out.println(10 / 3);    // 3  (parte decimal descartada)
System.out.println(10.0 / 3);  // 3.3333...

int a = 10, b = 3;
double resultado = (double) a / b;
System.out.println(resultado); // 3.3333...
\end{javacode}

% ─────────────────────────────────────────────────────────────────────────────
\section{Condicionais}

Condicionais permitem que o programa tome decisões.

\subsection{\texttt{if} e \texttt{else}}

\begin{javacode}
int temperatura = 30;

if (temperatura > 25) {
    System.out.println("Está quente!");
} else {
    System.out.println("Não está tão quente.");
}
\end{javacode}

\subsection{\texttt{else if}}

\begin{javacode}
int nota = 75;

if (nota >= 90) {
    System.out.println("A");
} else if (nota >= 80) {
    System.out.println("B");
} else if (nota >= 70) {
    System.out.println("C");
} else {
    System.out.println("Reprovado");
}
\end{javacode}

\subsection{Operadores de comparação}

\begin{center}
\begin{tabular}{c l}
\toprule
\textbf{Operador} & \textbf{Significado} \\
\midrule
\texttt{==} & Igual a       \\
\texttt{!=} & Diferente de  \\
\texttt{<}  & Menor que     \\
\texttt{>}  & Maior que     \\
\texttt{<=} & Menor ou igual \\
\texttt{>=} & Maior ou igual \\
\bottomrule
\end{tabular}
\end{center}

\subsection{Comparando Strings}

Para Strings, \textbf{não use \inlinecode{==}}. Use o método \inlinecode{.equals()}:

\begin{javacode}
String entrada = scanner.nextLine();

if (entrada.equals("sim")) {
    System.out.println("Você confirmou.");
} else {
    System.out.println("Você não confirmou.");
}
\end{javacode}

% ─────────────────────────────────────────────────────────────────────────────
\section{Resumo}

\begin{itemize}
  \item Todo programa Java tem a estrutura \inlinecode{public class ... \{ public static void main ... \}}.
  \item \inlinecode{System.out.println()} exibe texto e pula a linha; \inlinecode{System.out.print()} exibe sem pular.
  \item \inlinecode{Scanner} lê a entrada do usuário. Use \inlinecode{scanner.nextLine()} para ler texto.
  \item Os quatro tipos básicos são: \inlinecode{String}, \inlinecode{int}, \inlinecode{double} e \inlinecode{boolean}.
  \item Java realiza divisão inteira quando ambos os operandos são \inlinecode{int}.
  \item \inlinecode{if}, \inlinecode{else if} e \inlinecode{else} permitem que o programa tome decisões.
  \item Para comparar Strings, use \inlinecode{.equals()}, não \inlinecode{==}.
\end{itemize}

% ─────────────────────────────────────────────────────────────────────────────
\section{Exercícios}

\begin{enumerate}
  \item \textbf{Olá, usuário!} — Escreva um programa que pergunta o nome e a idade do usuário e exibe: \textit{Olá, [nome]! Você tem [idade] anos.}

  \item \textbf{Calculadora simples} — Escreva um programa que lê dois números inteiros e exibe a soma, diferença, produto e quociente (com resultado decimal).

  \item \textbf{Par ou ímpar} — Escreva um programa que lê um número inteiro e informa se é par ou ímpar.

  \item \textbf{Classificação de temperatura} — Lê uma temperatura em Celsius e exibe: \textit{Frio} (menor que 15), \textit{Agradável} (entre 15 e 25 inclusive) ou \textit{Quente} (maior que 25).

  \item \textbf{IMC} — Lê o peso (kg) e a altura (m) do usuário, calcula o IMC (\(peso / altura^2\)) e classifica: Abaixo do peso ($< 18{,}5$), Peso normal (18,5–24,9), Sobrepeso (25–29,9) ou Obesidade ($\geq 30$).
\end{enumerate}

\end{document}
